% 夜半用户手册 · 第 2 章：工程（中文）
% T-DOC-001 逐章交付。后续章节直接放 chapters/（依管理侧 #5278 对齐）。
% 事实来源：crates/yeban-app/ui/app.slint 菜单、host.rs 的工程/模板/启动器接线、
% startup_chooser.rs 的 NewProjectDetails、docs/manual/02-projects.md。

\chapter{工程}
\label{ch:projects}

\begin{quote}
九层之台，起于累土。
\end{quote}

出处：《道德经·第六十四章》

一个工程从空白容器和第一条轨道开始，先立骨架，再添内容。

这一章讲工程的整个生命周期：怎么建、怎么开、怎么存、怎么复用，以及夜半
「撤销就是版本」的独特模型。读完之后，你应该能放心地把一段工作存下来、
找回来，并回到任意一个过去的状态。

\section{新建、打开与保存}

\subsection{文件菜单里的五个入口}

「文件」菜单与工程直接相关的五项：

\begin{itemize}
  \item \textbf{新建空白工程}：得到一个只有主总线的空工程。
  \item \textbf{从模板新建...}：打开模板选择器，从工厂模板或你自己的模板起步。
  \item \textbf{打开 / 恢复工程...}：在文件对话框里选一个 \texttt{.yeban} 文件。
  \item \textbf{保存工程（Cmd/Ctrl+S）}：把当前工程写回原文件。
  \item \textbf{另存为（Cmd/Ctrl+Shift+S）}：换一个文件名或位置保存。
\end{itemize}

命令行也可以直接打开：\texttt{yeban-app --open <路径>}。

\screenshot{ch02-file-menu}{文件菜单展开，标出新建空白工程、从模板新建、打开/恢复、保存、另存为五项}

\subsection{启动工程选择器}

不指定工程启动时，会出现启动工程选择器。它把入口分成几类：\textbf{最近工程}、
\textbf{演示工程}、\textbf{工厂模板}、\textbf{用户模板}，以及停在「新建工程」一栏
的空白工程入口。选演示工程可以先逛一圈界面而不动你的作品；选最近工程则直接
回到上次的工作。

\screenshot{ch02-startup-chooser}{启动工程选择器，四个分类标签与「新建工程」一栏}

\subsection{保存是原子的}

保存采用原子写入：先把内容写到临时位置，确认完整后再替换目标文件。写一半
断电不会留下半个工程文件——要么是新版本，要么还是上一版。

\begin{figure}[htbp]
\centering
\begin{tikzpicture}[node distance=1.2cm, >=stealth,
    box/.style={draw, rounded corners, align=center, minimum width=2.4cm, minimum height=0.9cm}]
  \node[box] (a) {新建 / 从模板};
  \node[box, right of=a] (b) {编辑工程};
  \node[box, right of=b] (c) {保存（原子写入）};
  \node[box, below of=c] (d) {最近工程 / 再次打开};
  \node[box, below of=b] (e) {时光机回到旧版};
  \draw[->] (a) -- (b); \draw[->] (b) -- (c); \draw[->] (c) -- (d);
  \draw[->] (b) -- (e); \draw[->] (e) -- (b);
\end{tikzpicture}
\caption{工程生命周期：新建/从模板 $\to$ 编辑 $\to$ 保存 $\to$ 再打开；编辑中可随时回到旧版}
\label{fig:ch02-lifecycle}
\end{figure}

\section{工程文件与资产池}

工程文件的扩展名是 \texttt{.yeban}，它是唯一的工程格式。一个 \texttt{.yeban}
容器里装着：轨道、片段、音符、设备链，以及工程用到的素材（放在资产池里，
按内容哈希寻址）。

资产池的意义在于「工程自带素材」：把工程拷到另一台机器，用到的音频与 SFZ
样本会随容器同行，不必到处找文件。导出分轨、另存为模板时，资产池都按同一
来源取，与命令行和 MCP 的口径一致（见第 8 章）。

\begin{table}[htbp]
\centering
\caption{工程里几样东西各放在哪}
\label{tab:ch02-container}
\begin{tabular}{ll}
\hline
内容 & 在工程中的位置 \\
\hline
轨道 / 片段 / 音符 & 工程的编曲数据 \\
设备链 & 每条轨的设备序列 \\
提交图谱（history DAG） & 撤销链，随工程一起保存 \\
所用素材 & 资产池，按内容哈希寻址 \\
工程设置 & 速度 / 调号 / 拍号 / 采样率 \\
\hline
\end{tabular}
\end{table}

\section{模板：新建、另存为与工厂目录}

\subsection{从模板新建}

「文件 $\to$ 从模板新建...」打开模板选择器，列出两类模板：\textbf{工厂模板}
（随应用提供）与\textbf{用户模板}（你自己保存的）。选中一项即以其为骨架建出新
工程。列表很长时，选择器会提示被折叠的行数，不会悄悄丢掉条目。

\subsection{另存为模板}

「文件 $\to$ 另存为模板...」打开保存对话框：先填名字，确认后写入用户模板目录。
如果同名模板已存在，应用不会直接覆盖，而是停下来问你是否确认覆盖；取消则
什么也不写。用户模板目录在第一次使用时创建；目录不可用时，对话框会显示
错误原因，而不是静默失败。

\screenshot{ch02-template-dialog}{另存为模板对话框，显示命名输入、提示行与同名确认态}

\section{最近工程与恢复}

每次成功打开工程，路径都会被记进「最近工程」。它可以从启动工程选择器的
「最近工程」分类里找到并一键打开；列表里也提供移除入口，把不再需要的记录
删掉。最近工程、语言与布局共用同一份本机设置存储，所以改语言或界面布局
不会把最近工程挤掉。

「打开 / 恢复工程...」除了打开正常保存的工程，也用于恢复场景：选中容器文件
后，应用按其中记录的提交图谱与资产池还原现场。

\section{工程设置：采样率、速度、拍号}

在启动工程选择器的「新建工程」一栏，可以给空白工程设定几项基础参数，它们
会写进新工程：

\begin{itemize}
  \item \textbf{速度（BPM）}：可填范围内为 20--999；越界或非法输入会被收回到合法域。
  \item \textbf{调号}：可以不设（规范空工程的出厂态就是无调号），也可选升降号范围 $-7$ 到 $+7$。
  \item \textbf{拍号}：分子 1--32；分母只接受 1、2、4、8、16、32。
  \item \textbf{采样率}：工程的音频采样率（Hz）。
\end{itemize}

这几项决定后续换算：小节长度由拍号算出，时间码由工程拍号推出；拍号非法时
应用会如实报出原因（例如「拍号非法，无法计算片段长度」），而不是悄悄用
一个默认值顶上。

\screenshot{ch02-new-project-details}{新建工程一栏，显示速度、调号、拍号、采样率四个输入}

\section{撤销即版本：时光机}

夜半没有单独的「历史版本」功能：每一次编辑都是一次提交，撤销与重做只是在
这条提交链上移动。

\begin{itemize}
  \item \textbf{撤销}：Cmd/Ctrl+Z；\textbf{重做}：Cmd/Ctrl+Shift+Z。
  \item \textbf{编曲版本时光机（Undo Tree）}：「编辑 $\to$ 编曲版本时光机 (Undo Tree)...」；Cmd/Ctrl+Shift+H 也可打开。
\end{itemize}

时光机把整条提交链画成图谱：点任意一个版本节点，工程就回到那一版。如果
更早的版本太多、画布放不下，面板会如实报出没有画出的数量，而不是假装都
画了。提交图谱随工程一起保存，所以下次打开还能接着看。

\screenshot{ch02-undo-tree}{编曲版本时光机图谱，选中某个历史节点并高亮}

\section{帮助主题映射}

\begin{table}[htbp]
\centering
\caption{本章小节 $\to$ 应用内帮助锚点}
\label{tab:ch02-help}
\begin{tabular}{ll}
\hline
本章小节 & 帮助主题 \\
\hline
文件菜单里的五个入口 & \texttt{project-file-menu} \\
启动工程选择器 & \texttt{project-startup-chooser} \\
保存是原子的 & \texttt{project-atomic-save} \\
工程文件与资产池 & \texttt{project-container-assets} \\
从模板新建 & \texttt{project-template-new} \\
另存为模板 & \texttt{project-template-save} \\
最近工程与恢复 & \texttt{project-recent-recovery} \\
工程设置 & \texttt{project-settings} \\
撤销即版本：时光机 & \texttt{project-undo-tree} \\
\hline
\end{tabular}
\end{table}

\paragraph{本章学完能做什么} 独立完成工程的建、开、存、复用，并能用时光机回到任意历史版本；说清 \texttt{.yeban} 容器里装了什么。

\paragraph{动手练习} 用「从模板新建」建一个工程，改速度到 96 BPM、拍号 3/4；保存后另存为模板；故意做几次编辑，用 Cmd/Ctrl+Z 撤销再用时光机跳回最早的版本；最后关闭应用，从启动选择器的「最近工程」重新打开，核对设置还在。
